\documentclass[11pt]{amsart}

\usepackage{amsmath,amssymb,mathtools}
\usepackage{microtype}
\usepackage{booktabs}
\usepackage[hidelinks]{hyperref}

\hypersetup{
  pdftitle={Factorial-residue hypergraphs and covering congruences},
  pdfauthor={},
  pdfsubject={Weighted factorial-residue hypergraphs, covering congruences, and rough integers},
  pdfkeywords={factorials, edge-colored hypergraphs, rainbow edge covers, covering congruences, sieve methods, rough numbers}
}

\newtheorem{theorem}{Theorem}[section]
\newtheorem{proposition}[theorem]{Proposition}
\newtheorem{lemma}[theorem]{Lemma}
\newtheorem{corollary}[theorem]{Corollary}
\theoremstyle{definition}
\newtheorem{definition}[theorem]{Definition}
\theoremstyle{remark}
\newtheorem{remark}[theorem]{Remark}

\title[Factorial-residue hypergraphs]
{Factorial-residue hypergraphs and covering congruences}
\author{}
\date{}

\subjclass[2020]{11N36, 11B65, 11A41, 05C65}
\keywords{factorials, edge-colored hypergraphs, rainbow edge covers, covering congruences, sieve methods, rough numbers}

\begin{document}

\begin{abstract}
We study covering congruences for factorials using a weighted hypergraph
whose edges collect indices with equal factorial residues and whose
colors are primes. The minimum
weight $\kappa(L)$ of a rainbow cover equals the logarithm of the least
squarefree modulus of such a system. We prove
$\kappa(L)\leq(1-\delta)L\log L$ for an absolute $\delta>0$, even after
any $\lfloor L/100\rfloor$ primes are excluded. By counting permutations
of the remaining assignments, we obtain one modulus with
$((L+1)!)^{1/2-o(1)}$ distinct covering classes. This gives at least as
many $L$-rough integers in $(2L!,(L+1)!]$ whose differences from
$1!,\ldots,L!$ are all composite, uniformly in any prescribed reduced
residue class modulo $d\leq e^L$. We also prove
$\kappa(L)\geq(2/3-o(1))L^{1/3}\log L$. A resultant estimate bounds
the number of triples of equal factorial residues with small span. It
implies that covers whose consecutive triples all have span
$o(L^{1/4})$ have weight at least $(1/2-o(1))L\log L$.
The rough-integer result strengthens a quantitative solution of an
auxiliary problem of Erd\H{o}s; the integers obtained need not be prime.
\end{abstract}

\maketitle

\begin{center}
\small\textit{Working draft --- exposition and mathematical review in progress.}
\end{center}

\section{Introduction}\label{sec:intro}

In Problem A2 of \cite{Guy}, Erd\H{o}s asks whether there are infinitely
many primes $p$ such that $p-k!$ is composite whenever $1\leq k!<p$.
This is also listed as Erd\H{o}s Problem 1059 \cite{Bloom}. He also
suggested an auxiliary problem in which the integer itself need only
have large prime factors. In this paper, we study the covering
congruences used to construct such integers. When several factorials
have the same residue modulo a prime, one congruence makes that prime
divide each of the corresponding differences. We use these equalities
to reduce the product of the primes needed to cover all the factorials.

We first fix the notation for this covering problem. For $L\geq2$, let
$\mathcal H_L$ be the weighted edge-colored
multihypergraph on $V_L=\{2,\ldots,L\}$ whose nonempty edges are
\[
E_{q,a}=\{k\in V_L:k!\equiv a\pmod q\},
\qquad L<q<L!,\quad a\in(\mathbb Z/q\mathbb Z)^\times.
\]
Each edge has color $q$ and weight $\log q$. A \emph{rainbow cover} selects
at most one edge of each color and covers every vertex. Write $\kappa(L)$
for its minimum weight, with value $+\infty$ when no cover exists.
The Chinese remainder theorem (CRT) gives
\[
\kappa(L)=\log M_{\min}(L),
\]
where $M_{\min}(L)$ is the least squarefree modulus of a compatible
covering system using these primes and residues. We give the precise
definition and prove this equality in
Proposition~\ref{prop:crt-equivalence}. Our first result bounds the
weight of a cover.
\begin{theorem}\label{thm:covercost}
There are absolute constants $\delta>0$ and $L_0$ such that
\[
\kappa(L)\leq(1-\delta)L\log L
\]
for every integer $L\geq L_0$.
\end{theorem}
Assigning a different prime to each index gives a cover of weight
$(1+o(1))L\log L$. Theorem~\ref{thm:covercost} improves this bound by
pairing indices. The identity used in the proof is
\begin{equation}\label{eq:keyidentity}
(i+2)!-i!=i!\bigl(i^2+3i+1\bigr).
\end{equation}
A product argument and a dimension-one upper-bound sieve give a
positive proportion of pairs with primes $L<q\ll L^{2-\alpha}$.
We then use a conflict graph to choose disjoint pairs with distinct
primes. Each chosen prime covers both indices in its pair. We assign
individual primes of size $O(L\log L)$ to the remaining indices.

The same construction allows us to exclude any set of at most
$\lfloor L/100\rfloor$ primes; see Theorem~\ref{thm:robust}.
This will be useful when we prescribe an additional congruence for the
integers we construct. We also count the distinct residue classes
obtained by permuting the primes assigned to individual indices in
Theorem~\ref{thm:ensemble}.

We next consider lower bounds. Proposition~\ref{prop:fibresize} and
Corollary~\ref{cor:globallower} give
\[
\kappa(L)\geq(2/3-o(1))L^{1/3}\log L.
\]
There is a stronger bound when the indices in each fiber are close
together. Suppose every consecutive triple in every designated fiber
has span at most $cL^{1/4}$, where $c>0$ is fixed. Then
Corollary~\ref{cor:localcost} gives
\[
\frac{w(\mathcal C)}{\log((L+1)!)}
\geq\frac12-\frac{c^4}{192}-o(1).
\]
In particular, a full cover whose consecutive-triple spans are
$o(L^{1/4})$ costs at least $(1/2-o(1))L\log L$.
The order of $\kappa(L)$ is not determined by these bounds.

\subsection*{Erd\H{o}s's auxiliary problem}
The auxiliary problem of Erd\H{o}s asks for infinitely many pairs $(L,n)$
with $L!<n\leq(L+1)!$, every prime factor of $n$ exceeding $L$, and
$n-k!$ composite for $1\leq k\leq L$.
For $n>1$, write $P^-(n)$ and $P^+(n)$ for its least and greatest prime
factors. We call $n$ \emph{$L$-rough} when $P^-(n)>L$.

The covering construction gives the following quantitative form of the
auxiliary statement.
\begin{theorem}\label{thm:main}
There are absolute constants $\eta>0$ and $L_1$ such that, for every
integer $L\geq L_1$, at least
\[
\frac{((L+1)!)^\eta}{100\log L}
\]
integers $n\in(2L!,(L+1)!]$ are $L$-rough and have $n-k!$ composite for
all $1\leq k\leq L$.
\end{theorem}

By counting several covering classes, we obtain the following stronger
bound, also in prescribed arithmetic progressions.
\begin{theorem}\label{thm:abundance}
For every $\varepsilon>0$, there is $L_\varepsilon$ such that the following
holds for every integer $L\geq L_\varepsilon$. Put $X=(L+1)!$.
For every integer $1\leq d\leq e^L$ and every residue $b\pmod d$ with
$(b,d)=1$, at least $X^{1/2-\varepsilon}$ integers
$n\in(2L!,X]$ satisfy
\[
n\equiv b\pmod d,\qquad P^-(n)>L,
\qquad n-k!\text{ composite }(1\leq k\leq L).
\]
The threshold is independent of $d$ and $b$.
\end{theorem}

Our proof of Theorem~\ref{thm:abundance} uses the elementary estimate
\[
\#\{(i,j)\in\{1,\ldots,L\}^2:i!\equiv j!\pmod q\}
\leq6L^{3/2}\qquad(q>L),
\]
and Lemma~\ref{lem:permutationimage}, which bounds the number of
distinct images of permutations. Estimates of order $L^{3/2}$ for
equal factorial residues are classical; see
\cite[Theorem~1, with $\ell=1$]{GaraevLucaShparlinski}. A direct proof
with the displayed constant is included. We then count rough integers
in these classes by inclusion--exclusion.

Land \cite{Land} has posted a proposed solution of the prime problem
using a sieve for sparse sets of shifts and Fourier estimates. The
theorem stated there also implies existence in the auxiliary problem.
Our proofs are independent of that work. Here we study the size of the
covering modulus, the number of classes it gives, and restrictions on
the indices that one prime can cover.

Other work on factorial residues studies value sets and congruence
counts modulo a fixed prime; see
\cite{BanksLucaShparlinskiStichtenoth,GaraevLucaShparlinski,KlurmanMunsch,Hu}.
In our covering problem, the prime varies from one edge to another, but
only one residue can be chosen for each prime.

\subsection*{Outline of the paper}
Section~\ref{sec:hypergraph} defines the covers and proves the elementary
bounds used later. Sections~\ref{sec:quadratic}--\ref{sec:construction}
give the quadratic construction, including the choice of primes when
some are excluded. Section~\ref{sec:ensemble} counts covering classes,
and Section~\ref{sec:locality} bounds triples of equal factorial
residues with small span. In Section~\ref{sec:arithmetic}, we count
rough integers in the resulting progressions and prove
Theorems~\ref{thm:main} and \ref{thm:abundance}. We end with
computational examples and further questions.

The original quantitative auxiliary theorem and its proof using
quadratic congruences have an accompanying Lean formalization registered as
PALOMAR-2026-08-27-000003 v1 \cite{Palomar}. The standard dimension-one
upper-bound sieve is an explicit external hypothesis, matched to
\cite[Theorem~18.11(a)]{Koukoulopoulos}. That formalization does not
include the additional results here on excluded primes, counts of
covering classes, resultants, or prescribed progressions.

\section{Residue covers and covering congruences}\label{sec:hypergraph}

Let $L\geq2$ and put
\[
V_L=\{2,3,\ldots,L\}.
\]

\begin{definition}\label{def:fibres}
For a prime $L<q<L!$ and a nonzero residue $a\pmod q$, let $e_{q,a}$ be a
labeled hyperedge with incident vertex set
\[
E_{q,a}=\{k\in V_L:k!\equiv a\pmod q\}.
\]
After empty edges are omitted, these labeled edges form a weighted
edge-colored multihypergraph $\mathcal H_L$.  The edge $e_{q,a}$ has color
$q$ and weight $\log q$.
\end{definition}

Since $q>L$, every $k!$ with $k\in V_L$ is nonzero modulo $q$.  Therefore,
for each fixed $q$, the fibers $E_{q,a}$ partition $V_L$, with empty fibers
omitted.

\begin{definition}\label{def:cover}
A \emph{rainbow edge cover} of $\mathcal H_L$ is a finite collection
\[
\mathcal C=\{e_{q,a_q}:q\in Q\}
\]
with distinct primes $L<q<L!$, one residue $a_q\in(\mathbb Z/q\mathbb Z)^\times$
for each $q\in Q$, nonempty fibers $E_{q,a_q}$, and
\[
V_L\subseteq\bigcup_{q\in Q}E_{q,a_q}.
\]
Its weight is
\[
w(\mathcal C)=\sum_{q\in Q}\log q.
\]
We define
\[
\kappa(L)=\inf_{\mathcal C}w(\mathcal C),
\]
where the infimum is over all rainbow edge covers, and put
$\kappa(L)=+\infty$ if none exists.
\end{definition}

The rainbow condition is necessary because choosing two different fibers of
the same color would impose incompatible residues modulo the same prime.

\begin{proposition}\label{prop:criterion}
Let $2\leq i<j\leq L$ and let $q>L$ be prime.  Then
\[
\begin{aligned}
i!\equiv j!\pmod q
&\quad\Longleftrightarrow\quad
\prod_{r=i+1}^{j}r\equiv1\pmod q\\
&\quad\Longleftrightarrow\quad
q\mid\left(\prod_{r=i+1}^{j}r-1\right).
\end{aligned}
\]
\end{proposition}

\begin{proof}
We have
\[
j!=i!\prod_{r=i+1}^{j}r.
\]
Also, $q>L\geq i$, so $q\nmid i!$.  We may therefore cancel $i!$ modulo
$q$.  The three conditions are now equivalent.
\end{proof}

Assigning a different prime to each index gives the following bound.

\begin{proposition}\label{prop:baseline}
For every sufficiently large $L$,
\[
\kappa(L)\leq(L-1)\log(20L\log L)
=(1+o(1))L\log L.
\]
In particular, a rainbow edge cover exists, and the infimum in
Definition~\ref{def:cover} is a minimum.
\end{proposition}

\begin{proof}
The prime number theorem gives
\[
\pi(20L\log L)-\pi(L)>2L
\]
for all sufficiently large $L$.  Choose distinct primes
$q_k\in(L,20L\log L]$ for $k\in V_L$.  We also have
$20L\log L<L!$ for large $L$.  The edges $e_{q_k,k!}$ form a rainbow edge
cover, and their total weight is at most
\[
(L-1)\log(20L\log L).
\]
There are only finitely many primes between $L$ and $L!$, and each has only
finitely many fibers.  Hence there are only finitely many rainbow edge covers.
The nonempty finite set of their weights has a minimum.
\end{proof}

We now relate the weight of a cover to its modulus. For
sufficiently large $L$, let $M_{\min}(L)$ be the least squarefree product
$M$ of primes $L<q<L!$ for which there exists
$A\in(\mathbb Z/M\mathbb Z)^\times$ satisfying both of the following
conditions:
\begin{enumerate}
\item for every prime $q\mid M$, there is a $k\in V_L$ such that
$A\equiv k!\pmod q$;
\item for every $k\in V_L$, there is a prime $q\mid M$ such that
$A\equiv k!\pmod q$.
\end{enumerate}

\begin{proposition}[Covering congruences]
\label{prop:crt-equivalence}
For every sufficiently large $L$,
\[
\kappa(L)=\log M_{\min}(L).
\]
More precisely, a rainbow edge cover $\mathcal C$ gives a modulus
\[
M(\mathcal C)=\prod_{q\in Q}q
\]
with
\[
\log M(\mathcal C)=w(\mathcal C),
\]
and every pair $(M,A)$ in the definition of $M_{\min}(L)$ gives a rainbow
edge cover of weight $\log M$.
\end{proposition}

\begin{proof}
Let $\mathcal C=\{e_{q,a_q}:q\in Q\}$ be a rainbow edge cover.  The primes
in $Q$ are distinct, so the Chinese remainder theorem gives a unique residue class
$A\pmod {M(\mathcal C)}$ satisfying
\[
A\equiv a_q\pmod q\qquad(q\in Q).
\]
Every $a_q$ is nonzero, so $(A,M(\mathcal C))=1$.  If $k\in E_{q,a_q}$,
then $A\equiv k!\pmod q$.  Thus $(M(\mathcal C),A)$ satisfies the defining
conditions for $M_{\min}(L)$.  Also,
\[
\log M(\mathcal C)=\sum_{q\in Q}\log q
=w(\mathcal C).
\]

Conversely, let $(M,A)$ satisfy the defining conditions for $M_{\min}(L)$.
For each prime $q\mid M$, choose the edge $e_{q,A\bmod q}$.  The residue is
nonzero because $(A,M)=1$, and the second condition says that these edges
cover $V_L$.  They use at most one edge of each color and have total weight
\[
\sum_{q\mid M}\log q=\log M.
\]
Taking the two minima proves the proposition.
\end{proof}

\subsection{Finite covers and a lower bound}
Let $m\geq1$, let $\mathbf h=(h_1,\ldots,h_m)$ be distinct nonzero
integers, and let
$\mathcal P$ be a finite set of primes, none dividing $\prod_i h_i$.
For $q\in\mathcal P$, use the nonempty fibers
\[
E_{q,a}(\mathbf h)=\{i:h_i\equiv a\pmod q\}
\]
with weight $\log q$. Define $\kappa_{\mathcal P}(\mathbf h)$ by the same
definition, with value $+\infty$ if no cover exists.
The proof of Proposition~\ref{prop:crt-equivalence} applies here as
well. Choosing at most one nonzero residue for each prime gives a
compatible system, and its weight is $\log M$. The minimum depends on
the specified prime set $\mathcal P$.

Equivalently, the minimum is given by the integer program
\[
\min \sum_{q,a}x_{q,a}\log q,
\qquad
\sum_{q,a:i\in E_{q,a}}x_{q,a}\geq1,
\qquad
\sum_a x_{q,a}\leq1,
\qquad x_{q,a}\in\{0,1\}.
\]
The sums run over nonempty fibers. Allowing real values
$x_{q,a}\geq0$ gives the corresponding linear programming relaxation.
Its minimum is a lower bound for the minimum above.

\begin{proposition}[A lower bound for the cover weight]
\label{prop:dual}
Let $y_1,\ldots,y_m$ be nonnegative real numbers, and set
\[
z_q=\max\left\{0,\ \max_a\sum_{i\in E_{q,a}(\mathbf h)}y_i-\log q\right\}.
\]
Then
\[
\kappa_{\mathcal P}(\mathbf h)\geq\sum_{i=1}^m y_i-\sum_{q\in\mathcal P}z_q.
\]
The same inequality holds for every feasible fractional solution above.
\end{proposition}
\begin{proof}
For a feasible vector $x$, nonnegativity, the vertex constraints, and then
the color constraints give
\[
\begin{aligned}
\sum_i y_i
&\leq\sum_{q,a}x_{q,a}\sum_{i\in E_{q,a}}y_i\\
&\leq\sum_{q,a}x_{q,a}\log q+\sum_q z_q\sum_a x_{q,a}\\
&\leq\sum_{q,a}x_{q,a}\log q+\sum_qz_q.
\end{aligned}
\]
Rearranging proves the assertion for both integer and fractional covers.
\end{proof}

\begin{proposition}[Consecutive integer shifts]\label{prop:intervalshifts}
For $\mathbf h=(1,2,\ldots,m)$ and an admissible set $\mathcal P$ of at
least $m$ primes all exceeding $m$, let $q_1<\cdots<q_m$ be its $m$
smallest members. Then
\[
\kappa_{\mathcal P}(\mathbf h)=\sum_{j=1}^m\log q_j.
\]
If the set $\mathcal P$ contains the first $m$ primes exceeding $m$,
then this minimum is asymptotic to $m\log m$.
\end{proposition}
\begin{proof}
Reduction modulo $q>m$ is injective on $\{1,\ldots,m\}$. Every fiber is a
singleton, so at least $m$ different primes are necessary, and the $m$
smallest suffice. In the last assertion the lower bound is $m\log m$.
The prime number theorem places the $m$th prime exceeding $m$ below
$2m\log m$ for sufficiently large $m$, giving the matching upper bound
$m\log(2m\log m)=(1+o(1))m\log m$.
\end{proof}

\subsection{Choosing compatible edges}
To choose edges with distinct primes, let $\mathcal F$ be a finite family
of triples
\[
F=(q,a,B_F),
\qquad
\varnothing\neq B_F\subseteq E_{q,a}.
\]
The set $B_F$ is the \emph{designated subset} of $F$.  Define the auxiliary
conflict graph $G_{\mathcal F}$ on vertex set $\mathcal F$ by joining
distinct triples $F=(q,a,B_F)$ and $F'=(q',a',B_{F'})$ whenever $q=q'$ or
\[
B_F\cap B_{F'}\neq\varnothing.
\]

\begin{proposition}[An independent set of large weight]
\label{prop:conflict-selection}
Suppose that the conflict graph of $\mathcal F$ has maximum degree at most
$D$.  If $w(F)\geq0$ for every $F\in\mathcal F$, then there is a subfamily
$\mathcal J\subseteq\mathcal F$ with distinct prime colors and pairwise
disjoint designated subsets such that
\[
\sum_{F\in\mathcal J}w(F)
\geq\frac1{D+1}\sum_{F\in\mathcal F}w(F).
\]
In the special case $w(F)=1$ for every $F$, one may choose $\mathcal J$ so
that
\[
|\mathcal J|\geq\frac{|\mathcal F|}{D+1}.
\]
\end{proposition}

\begin{proof}
A finite graph of maximum degree at most $D$ has a proper coloring with at
most $D+1$ colors: order the vertices arbitrarily and color them greedily.
The color classes are independent sets and partition $\mathcal F$.  One of
them has weight at least the average of their weights.  Taking this color
class as $\mathcal J$ proves the weighted inequality.  Taking unit weights
gives the last assertion.  Independence gives distinct prime colors and
pairwise disjoint designated subsets by the definition of the conflict graph.
\end{proof}

Covering the indices in $B_F$ individually has leading weight
$|B_F|\log L$. A single prime saves a fixed proportion of this weight
only when $q\leq L^{|B_F|-\varepsilon}$ for some $\varepsilon>0$.
In the quadratic construction $|B_F|=2$, so we seek primes in the range
$q\ll L^{2-\alpha}$.


\begin{proposition}[Completing a cover]\label{prop:completion}
In the finite shift system above, suppose a candidate family $\mathcal F$
has designated nonempty sets $B_F\subseteq E_{q_F,a_F}$, and its conflict
graph has maximum degree at most $\Delta$. Suppose there are at least $m$
additional admissible primes, all at most $T$, none occurring as a
candidate color. Define
\[
\sigma(F)=\bigl(|B_F|\log T-\log q_F\bigr)_+,
\qquad x_+=\max\{x,0\}.
\]
Then the enlarged prime set admits a rainbow cover of weight at most
\[
m\log T-\frac1{\Delta+1}\sum_{F\in\mathcal F}\sigma(F).
\]
\end{proposition}
\begin{proof}
Discard candidates with zero saving and apply
Proposition~\ref{prop:conflict-selection} with weights $\sigma(F)$.
The selected family $\mathcal J$ has disjoint designated sets and
saves at least $(\Delta+1)^{-1}\sum_F\sigma(F)$.
Assign a different additional prime to each remaining index. If
$v=\sum_{F\in\mathcal J}|B_F|$, the total weight is at most
\[
\sum_{F\in\mathcal J}\log q_F+(m-v)\log T
=m\log T-\sum_{F\in\mathcal J}\sigma(F).
\]
All added primes are different from the selected colors, so the cover is
compatible.
\end{proof}

\section{Prime divisors of a quadratic}\label{sec:quadratic}

We first count the indices for which $i^2+3i+1$ has a sufficiently
large prime divisor. Throughout this section, put
\[
f(t)=t^2+3t+1.
\]

\subsection{Roots modulo prime powers}

We also let $\chi_5$ be the primitive quadratic character modulo $5$, and we
extend it by setting $\chi_5(5)=0$.  We will use the prime number theorem, the
Siegel--Walfisz theorem for the fixed modulus $5$, and Mertens's estimates in
the forms stated in
\cite[Theorems 3.4, 8.1 and 12.1]{Koukoulopoulos}.  Let
\[
\Theta_{\chi_5}(x)=\sum_{p\leq x}\chi_5(p)\log p.
\]
For $a\in(\mathbb Z/5\mathbb Z)^\times$, the fixed-modulus Siegel--Walfisz
estimate for $\pi(x;5,a)$ and partial summation give
$\vartheta(x;5,a)=x/4+O(xe^{-c\sqrt{\log x}})$ for some $c>0$.  Since the
sum of $\chi_5(a)$ over these four classes is zero, summing with weights
$\chi_5(a)$ gives
$\Theta_{\chi_5}(x)\ll xe^{-c\sqrt{\log x}}$.
Partial summation therefore gives
\begin{equation}\label{eq:chi}
\begin{aligned}
\sum_{p\leq x}\chi_5(p)\frac{\log p}{p-1}
&=\frac{\Theta_{\chi_5}(x)}{x-1}
 +\int_2^x\frac{\Theta_{\chi_5}(t)}{(t-1)^2}\,dt+O(1)\\
&=O(1).
\end{aligned}
\end{equation}
Indeed, the integrand is $O(e^{-c\sqrt{\log t}}/t)$, which is integrable on
$[2,\infty)$.  The same argument, now using the weight $1/(t\log t)$,
gives
\[
\sum_{p\leq x}\frac{\chi_5(p)}p
=\frac{\Theta_{\chi_5}(x)}{x\log x}
 +\int_2^x\Theta_{\chi_5}(t)
 \frac{\log t+1}{t^2(\log t)^2}\,dt+O(1)
=O(1).
\]
For $p\neq5$, we have
$\mathbf 1_{\chi_5(p)=1}=(1+\chi_5(p))/2$.  Mertens's estimate now gives
\begin{equation}\label{eq:split}
\sum_{\substack{p\leq x\\ \chi_5(p)=1}}\frac1p
=\frac12\log\log x+O(1).
\end{equation}

The identity
\begin{equation}\label{eq:discriminant}
(2t+3)^2=4f(t)+5
\end{equation}
shows that, for $p\neq2,5$, the number of roots of $f$ modulo $p$ is
$1+\left(\frac5p\right)$.  Since $5\equiv1\pmod4$, quadratic reciprocity
gives
\[
\left(\frac5p\right)=\left(\frac p5\right)=\chi_5(p).
\]
Thus,
\begin{equation}\label{eq:rho}
\rho(p)=1+\chi_5(p)\in\{0,2\}.
\end{equation}
The derivative is $f'(t)=2t+3$.  If a root of $f$ modulo $p$ were not simple,
then $f(t)\equiv f'(t)\equiv0\pmod p$, and \eqref{eq:discriminant} would give
$p\mid5$.  Hence, every root under consideration is simple.  Hensel's lemma
then gives exactly one lift of each root to every modulus $p^a$.  There is no
root modulo $2$.  Modulo $5$ there is one root, but it does not lift to a root
modulo $25$, because
\begin{equation}\label{eq:five}
f(1+5t)=5(1+5t+5t^2).
\end{equation}

\subsection{Large prime factors}

\begin{lemma}\label{lem:largeprime}
For all sufficiently large $N$,
\[
\#\left\{1\leq i\leq N:P^+(f(i))>\frac18 i\log i\right\}\geq\frac N5.
\]
\end{lemma}

\begin{proof}
We begin with the indices for which every prime factor of $f(i)$ is at
most $i$. Put
\[
E_N=\left\{i:\frac{N}{\log N}<i\leq N,\ P^+(f(i))\leq i\right\},
\qquad
Q_N=\prod_{i\in E_N}f(i).
\]
Every prime divisor of $Q_N$ is at most $N$.  For $p\neq2,5$, the root count
modulo prime powers gives
\begin{align*}
\operatorname{ord}_p\left(\prod_{i\leq N}f(i)\right)
&=\sum_{a\geq1}\#\{i\leq N:p^a\mid f(i)\}.
\end{align*}
For every $a\geq1$, there are $\rho(p)$ root classes modulo $p^a$, and hence
\[
\#\{i\leq N:p^a\mid f(i)\}
\leq \frac{N\rho(p)}{p^a}+\rho(p).
\]
Also, $p^a\mid f(i)$ is possible only when $p^a\leq N^2+3N+1$.
Consequently, only $O(\log N/\log p)$ values of $a$ occur, and summing the
geometric series gives
\[
\operatorname{ord}_p\left(\prod_{i\leq N}f(i)\right)
\leq
\frac{N\rho(p)}{p-1}
+O\left(\rho(p)\frac{\log N}{\log p}\right).
\]
The prime $2$ contributes nothing.  Equation \eqref{eq:five} shows that
$5\mid f(i)$ precisely when $i\equiv1\pmod5$, and that in this case
$5^2\nmid f(i)$.  Thus, the contribution of $p=5$ to the logarithm below is
$O(N)$.  Moreover,
\[
\log N\sum_{p\leq N}\rho(p)\ll \log N\,\pi(N)\ll N.
\]
Therefore, Mertens's estimate together with \eqref{eq:chi} gives
\begin{align*}
\log Q_N
&\leq
N\sum_{p\leq N}\frac{(1+\chi_5(p))\log p}{p-1}
+O\left(\log N\sum_{p\leq N}\rho(p)\right)+O(N)\\
&=N\log N+O(N).
\end{align*}
For the corresponding lower bound, uniformly for
$N/\log N<i\leq N$, we have
\[
\log f(i)=2\log N+O(\log\log N).
\]
It follows that
\[
\log Q_N
\geq |E_N|\bigl(2\log N-O(\log\log N)\bigr).
\]
Comparing this inequality with the upper estimate for $\log Q_N$, we obtain
\[
|E_N|\leq\left(\frac12+o(1)\right)N.
\]
The interval $(N/\log N,N]$ contains $N-o(N)$ integers.  Thus, at least
$(1/2-o(1))N$ integers in this interval satisfy $P^+(f(i))>i$.  Put
\[
\mathcal S_N=
\left\{i:\frac{N}{\log N}<i\leq N,\ P^+(f(i))>i\right\}.
\]
We have proved that $|\mathcal S_N|\geq(1/2-o(1))N$.  It remains to
strengthen the lower bound for these prime factors.

Discard from $\mathcal S_N$ the indices for which
\[
P^+(f(i))\leq\frac18N\log N.
\]
For every discarded index, the prime $q=P^+(f(i))$ exceeds $i$.  Hence, the
relevant indices lie in $[1,q-1]$ and are distinct modulo $q$.  Since $f$
has at most two roots modulo $q$, a fixed prime $q$ can occur for at most two
indices.  Therefore,
the number of discarded indices is at most
\[
2\pi\left(\frac18N\log N\right)
=\left(\frac14+o(1)\right)N.
\]
At least $(1/2-o(1))N$ indices were available, and at most
$(1/4+o(1))N$ were discarded. Thus at least $(1/4-o(1))N$ remain. For
each remaining index,
\[
P^+(f(i))>\frac18N\log N\geq\frac18i\log i.
\]
Since $1/4-o(1)>1/5$ for all sufficiently large $N$, at least $N/5$ indices
satisfy the required inequality.  This proves the lemma.
\end{proof}

\section{A sieve estimate for the complementary factor}\label{sec:sieve}

A prime used to cover a pair must also satisfy the upper bound needed
for the modulus. Write $f(i)=mq$, where $q>L$ is prime. A lower bound
for $m$ gives an upper bound for $q$. We therefore estimate the number
of indices for which $m$ is small.

\subsection{The polynomials to be sieved}

For $0<\alpha\leq1/2$, let
\[
E_\alpha(L)
=
\#\left\{1\leq i\leq L:
 f(i)=mq\text{ for some }m\leq L^\alpha
 \text{ and prime }q>L
\right\}.
\]
Thus, $E_\alpha(L)$ counts the indices for which a prime factor $q>L$ has a
complementary factor $m\leq L^\alpha$. The next lemma bounds this number
in terms of $\alpha$.

\begin{lemma}\label{lem:smallcofactor}
There is an absolute constant $C_0\geq1$ such that, uniformly for $0<\alpha\leq1/2$ and all sufficiently large $L$,
\[
E_\alpha(L)
\leq
C_0\left(
\alpha L+\frac{L}{\log L}
+L^{(1+\alpha)/2}(\log L)^2
\right).
\]
The lower threshold for $L$ may be chosen independently of $\alpha$.
\end{lemma}

\begin{proof}
We begin by counting the roots of $f$ modulo $m$.  For $m\geq1$, write
\[
R(m)=\#\{a\pmod m:f(a)\equiv0\pmod m\}.
\]
The prime-power values from the previous section are
\[
R(2^e)=0\quad(e\geq1),
\qquad
R(5)=1,
\qquad
R(5^e)=0\quad(e\geq2),
\]
and, for $p\neq2,5$,
\[
R(p^e)=1+\chi_5(p)\quad(e\geq1).
\]
The Chinese remainder theorem therefore gives
\begin{equation}\label{eq:R}
R(m)=
\begin{cases}
2^{\omega(m/(m,5))},
&\begin{array}{l}
2\nmid m,\ 25\nmid m,\ \text{and}\\
\chi_5(p)=1\text{ for every }p\mid m,\ p\neq5,
\end{array}\\[3mm]
0,&\text{otherwise}.
\end{cases}
\end{equation}

Now fix $m\leq L^\alpha$ and a root $a\pmod m$, represented by
$0\leq a<m$.  Every integer $i\equiv a\pmod m$ can be written as
$i=a+mt$.  Accordingly, let
\[
i=a+mt,
\qquad
T_{a,m}=\{t\in\mathbb Z:1\leq a+mt\leq L\},
\qquad
H=|T_{a,m}|.
\]
The length of this interval of $t$ values is
\[
H=\frac Lm+O(1),
\qquad
H\asymp\frac Lm,
\qquad
H\geq L^{1/2}-1.
\]
Since $m\mid f(a)$, we can divide $f(a+mt)$ by $m$ and obtain the integral
polynomial
\[
G_{a,m}(t)=\frac{f(a+mt)}m
=mt^2+(2a+3)t+\frac{f(a)}m.
\]
This polynomial has discriminant $5$.  It is also primitive.  Indeed, if a
prime divided all three coefficients, then its square would divide
$(2a+3)^2-4f(a)=5$, which is impossible.

If $\ell\nmid10m$ is prime, then the affine map $t\mapsto a+mt$ is a
bijection modulo $\ell$.  Consequently,
\[
\#\{t\pmod\ell:G_{a,m}(t)\equiv0\pmod\ell\}
=1+\chi_5(\ell).
\]
Let
\[
\mathcal P_0=\{\ell:\ell\text{ is prime and }\ell\nmid10\},
\qquad
\mathcal P_m=\{\ell\in\mathcal P_0:\ell\nmid m\},
\]
and extend $\rho$ multiplicatively to squarefree integers supported on
$\mathcal P_m$, with $\rho(d)=0$ otherwise.  For supported squarefree $d$,
counting the root classes inside $T_{a,m}$ gives
\begin{equation}\label{eq:localcount}
\#\{t\in T_{a,m}:d\mid G_{a,m}(t)\}
=H\frac{\rho(d)}d+r_d,
\qquad
|r_d|\leq\rho(d).
\end{equation}

\subsection{Application of the fundamental lemma}

We apply the fundamental lemma with constants independent of $m$ and
$a$. Define the nonnegative weights
\[
a_n=\#\{t\in T_{a,m}:G_{a,m}(t)=n\},
\qquad X_{a,m}=\sum_n a_n=H.
\]
Thus, repeated values of $G_{a,m}$ are counted with their multiplicities.
For $d$ supported on $\mathcal P_m$, equation \eqref{eq:localcount} says
\[
\sum_{\substack{n\geq1\\d\mid n}}a_n
=H\frac{\rho(d)}d+r_d.
\]
This is precisely Axiom~1 of \cite[pp.~185--186]{Koukoulopoulos}, with local
density $\nu(d)=\rho(d)$ and remainder $r_d$.  We also have
$0\leq\rho(\ell)<\ell$ for every $\ell\in\mathcal P_m$.

To check the dimension condition, Mertens's estimate and
\eqref{eq:chi} give, uniformly for $w\geq2$,
\begin{equation}\label{eq:weightedrho}
\sum_{\substack{\ell\leq w\\ \ell\nmid10}}
\frac{\rho(\ell)\log\ell}{\ell}
=\sum_{\ell\leq w}\frac{\log\ell}{\ell}
+\sum_{\ell\leq w}\chi_5(\ell)\frac{\log\ell}{\ell}
+O(1)
=\log w+O(1).
\end{equation}
Thus, every finite truncation of the base prime set $\mathcal P_0$ satisfies
Axiom~$2'$ of \cite[p.~188]{Koukoulopoulos} with $\kappa=1$, uniformly in
the truncation point.  The additional local condition in that axiom holds,
for example, with $k=2$ and $\varepsilon=1/2$, since
$\rho(\ell)\in\{0,2\}$ and every prime in $\mathcal P_0$ for which
$\rho(\ell)=2$ is at least $11$.  The cited source therefore gives Axiom~2
for every finite truncation of $\mathcal P_0$, with the same absolute
constant $C$.

To obtain Axiom~2 for $\mathcal P_m$, observe that deleting primes
can only decrease its defining product. If
$3/2\leq y_1\leq y_2$, then
\[
\prod_{\substack{\ell\in\mathcal P_m\cap(y_1,y_2]}}
\left(1-\frac{\rho(\ell)}\ell\right)^{-1}
\leq
\prod_{\substack{\ell\in\mathcal P_0\cap(y_1,y_2]}}
\left(1-\frac{\rho(\ell)}\ell\right)^{-1}.
\]
The same inequality holds for each finite truncation. Axiom~2 for
$\mathcal P_0$ therefore gives Axiom~2 for $\mathcal P_m$, with the
same absolute constants. In particular, these constants are independent
of $m$ and $a$.

Set $z=H^{1/20}$ and
\[
\mathcal Q_{m,z}=\{\ell\in\mathcal P_m:\ell\leq z\},
\qquad
P_m(z)=\prod_{\ell\in\mathcal Q_{m,z}}\ell.
\]
For sufficiently large $L$, we have $z<L$.  Therefore, every prime value
$G_{a,m}(t)>L$ is counted by the sifted set
$\{t\in T_{a,m}:(G_{a,m}(t),P_m(z))=1\}$.  Let $S_{a,m}$ denote the number
of $t\in T_{a,m}$ for which $G_{a,m}(t)$ is a prime exceeding $L$.  If
$\mathcal Q_{m,z}$ is empty, the estimate below is immediate.  Otherwise,
let $y$ be its largest prime.  We apply
\cite[Theorem 18.11(a), p.~190]{Koukoulopoulos} to the finite prime set
$\mathcal Q_{m,z}$ with $u=10$.  In the notation of that theorem, its main
term is
\[
\left(1+O_{\kappa,C}(u^{-u/2})\right)H V_m(z),
\]
and its remainder is bounded by
\[
O\left(\sum_{\substack{d\leq y^u\\d\mid P_m(z)}}|r_d|\right).
\]
Since $y\leq z$, we have $y^u\leq z^{10}=H^{1/2}$.  Taking an upper bound
and enlarging the remainder sum, we obtain
\begin{equation}\label{eq:Sam}
S_{a,m}
\ll
HV_m(z)
+\sum_{\substack{d\leq H^{1/2}\\ d\mid P_m(z)}}|r_d|
\leq HV_m(z)+\sum_{d\leq H^{1/2}}\rho(d),
\end{equation}
where
\[
V_m(z)=
\prod_{\substack{\ell\leq z\\ \ell\nmid10m}}
\left(1-\frac{\rho(\ell)}\ell\right).
\]
For the remainder, $\rho(d)\leq\tau(d)$ gives
\[
\sum_{d\leq H^{1/2}}\rho(d)
\leq\sum_{ab\leq H^{1/2}}1
\ll H^{1/2}\log H.
\]
It remains to estimate the sieve product.  First put
\[
V_0(z)=\prod_{\substack{\ell\leq z\\ \ell\nmid10}}
\left(1-\frac{\rho(\ell)}\ell\right).
\]
By \eqref{eq:split} and the convergence of $\sum_p p^{-2}$, we have
\[
\log V_0(z)
=\sum_{\substack{\ell\leq z\\ \chi_5(\ell)=1}}
\log\left(1-\frac2\ell\right)
=-2\sum_{\substack{\ell\leq z\\ \chi_5(\ell)=1}}\frac1\ell+O(1)
=-\log\log z+O(1).
\]
To account for the primes dividing $m$, define
\[
h(m)=
\prod_{\substack{p\mid m,\ p\neq5\\ \chi_5(p)=1}}
\left(1-\frac2p\right)^{-1}.
\]
Then, omitting from $V_0(z)$ the primes that divide $m$ gives
\[
V_m(z)
=V_0(z)
\prod_{\substack{\ell\mid m,\ \ell\leq z\\
\ell\neq5,\ \chi_5(\ell)=1}}
\left(1-\frac2\ell\right)^{-1}
\ll\frac{h(m)}{\log z}
\ll\frac{h(m)}{\log L}
\]
uniformly for $m\leq L^\alpha$.  The last comparison follows from
$H\geq L^{1/2}-1$, which gives $\log z\asymp\log L$.  Substituting this into
\eqref{eq:Sam}, we obtain
\begin{equation}\label{eq:Sam2}
S_{a,m}
\ll
\frac{Hh(m)}{\log L}+H^{1/2}\log L.
\end{equation}

\subsection{Summation over complementary factors}

If $i$ is counted by $E_\alpha(L)$,
then for at least one representation $f(i)=mq$, the residue class
$a\equiv i\pmod m$ is one of the $R(m)$ roots and $G_{a,m}(t)=q$ for the
corresponding $t$. Summing over $m$ and its root classes gives an upper
bound, since an index may have more than one such representation.
Since $m\leq L^\alpha\leq L^{1/2}$, we also have
\[
H\leq\frac Lm+1\leq\frac{2L}{m}.
\]
Put $Y:=L^\alpha$.  Equation \eqref{eq:Sam2} now gives
\begin{equation}\label{eq:Esum}
E_\alpha(L)
\ll
\frac{L}{\log L}
\sum_{m\leq Y}\frac{R(m)h(m)}m
+L^{1/2}\log L
\sum_{m\leq Y}\frac{R(m)}{\sqrt m}.
\end{equation}
The two sums on the right are estimated separately.  For the first, define
$g(m)=R(m)h(m)$.  Equation \eqref{eq:R} and the definition of $h$ show that
$g$ is nonnegative and multiplicative, and that it vanishes outside the
admissible support.  If $p\neq5$ and $\chi_5(p)=1$, then
\[
1+\sum_{e\geq1}\frac{g(p^e)}{p^e}
=1+\frac{2p}{(p-1)(p-2)}
=1+\frac2p+O(p^{-2}).
\]
The prime $5$ contributes the bounded local factor $1+1/5$, while inert
primes and $2$ contribute $1$.  Therefore, positivity, the Euler product,
and \eqref{eq:split} give, for $Y\geq1$,
\begin{align*}
\sum_{m\leq Y}\frac{R(m)h(m)}m
&\leq
\prod_{p\leq Y}\left(1+\sum_{e\geq1}\frac{g(p^e)}{p^e}\right)\\
&\ll
\exp\left(
2\sum_{\substack{p\leq Y\\ \chi_5(p)=1}}\frac1p+O(1)
\right)
\ll\log(2Y).
\end{align*}
For the second sum, we use $R(m)\leq\tau(m)$.  Thus,
\begin{align*}
\sum_{m\leq Y}\frac{R(m)}{\sqrt m}
&\leq
\sum_{ab\leq Y}\frac1{\sqrt{ab}}\\
&\ll
\sum_{a\leq Y}\frac1{\sqrt a}\sqrt{\frac Ya}
\ll\sqrt Y\log(2Y).
\end{align*}
Finally,
\[
\log(2Y)\leq\alpha\log L+O(1),
\]
uniformly for $0<\alpha\leq1/2$.  Substitution of these bounds into
\eqref{eq:Esum} gives
\[
E_\alpha(L)
\ll \alpha L+\frac{L}{\log L}
+L^{(1+\alpha)/2}(\log L)^2,
\]
with an absolute implied constant and an $L$-threshold independent of
$\alpha$.  This proves the lemma.
\end{proof}

\section{Covers from quadratic congruences}\label{sec:construction}

We now combine Lemmas~\ref{lem:largeprime} and \ref{lem:smallcofactor}.
We choose disjoint pairs with distinct primes and assign individual
primes to the remaining indices.

\begin{proposition}\label{prop:quadratic-cover}
Let $C_0$ be as in Lemma~\ref{lem:smallcofactor}, and set
\[
\alpha_0=\frac1{1000C_0},
\qquad
\delta_0=\frac{\alpha_0}{200},
\qquad
\eta_0=\frac{\delta_0}{2}.
\]
For every sufficiently large $L$, there exist a set
$I\subseteq\{2,\ldots,L-2\}$, distinct primes $q_i$ for $i\in I$, and
distinct primes $r_k$, indexed by
\[
U=\{2,\ldots,L\}\setminus\bigcup_{i\in I}\{i,i+2\}.
\]
Put
\[
M=\prod_{i\in I}q_i\prod_{k\in U}r_k.
\]
These objects satisfy the following conditions:
\begin{enumerate}
\item The pairs $\{i,i+2\}$, $i\in I$, are pairwise disjoint, and
$|I|\geq L/100$.
\item For every $i\in I$,
\[
L<q_i<2L^{2-\alpha_0},
\qquad
q_i\mid f(i).
\]
\item For every $k\in U$,
\[
L<r_k\leq20L\log L.
\]
\item The primes $q_i$, $i\in I$, and $r_k$, $k\in U$, are pairwise
distinct.
\item The modulus $M$ satisfies
\[
\log M\leq(1-\delta_0)L\log L,
\qquad
M\leq ((L+1)!)^{1-\eta_0}.
\]
\end{enumerate}
Moreover, every prime divisor of $M$ is less than $L!$.
\end{proposition}

\begin{proof}
The choice of $\alpha_0$ in the proposition and
Lemma~\ref{lem:smallcofactor} give
\[
C_0\alpha_0L=\frac L{1000}.
\]
Since $\alpha_0<1$ is fixed, we also have
\[
C_0\left(\frac L{\log L}
+L^{(1+\alpha_0)/2}(\log L)^2\right)=o(L).
\]
Consequently,
\[
E_{\alpha_0}(L)\leq\frac L{500}
\]
for sufficiently large $L$.  On the other hand,
Lemma~\ref{lem:largeprime} supplies at least $L/5$ indices $i\leq L$ for
which
\[
q_i:=P^+(f(i))>\frac18i\log i.
\]
We discard the indices $i<L/10$, the two endpoints $L-1,L$, and those
counted by $E_{\alpha_0}(L)$.  The number remaining is at least
\[
\frac L5-\frac L{10}-\frac L{500}-2
=\frac{49L}{500}-2
\geq\frac{9L}{100}
\]
for sufficiently large $L$.  For every remaining index,
\[
q_i>\frac18 i\log i
\geq \frac{L}{80}\log\left(\frac{L}{10}\right)>L
\]
once $L$ is sufficiently large.  Write $f(i)=m_iq_i$.  Then
$m_i>L^{\alpha_0}$, since otherwise this index would be counted by
$E_{\alpha_0}(L)$.  Since $f(i)<2L^2$ for $i\leq L$, it follows that
\[
q_i<2L^{2-\alpha_0}.
\]
The identity \eqref{eq:keyidentity} now gives
\[
i!\equiv(i+2)!\pmod{q_i}.
\]

For each remaining index $i$, let
\[
F_i=(q_i,i!,B_i),
\qquad B_i=\{i,i+2\}.
\]
The identity above shows that $B_i\subseteq E_{q_i,i!}$.  We apply
Proposition~\ref{prop:conflict-selection}.  Two designated pairs
$\{i,i+2\}$ and $\{j,j+2\}$ intersect only when $j=i$ or $j=i\pm2$.
Thus, a fixed vertex of the auxiliary conflict graph has at most two
neighbors arising from an overlap of designated pairs.

A fixed prime $q>L$ occurs for at most two starting indices, because all
starts are less than $q$ and $f$ has at most two roots modulo $q$.  Hence a
fixed vertex has at most one additional neighbor of the same prime color.
The conflict graph therefore has maximum degree at most $3$.
Proposition~\ref{prop:conflict-selection}, with unit weights, gives at least
\[
\frac14\cdot\frac{9L}{100}=\frac{9L}{400}
\]
vertices whose associated edges have distinct colors and whose designated
pairs are pairwise disjoint.  For
sufficiently large $L$, we may therefore choose a subfamily indexed by a set
$I$ with
\[
|I|=\left\lceil\frac L{100}\right\rceil.
\]

The set of remaining indices has cardinality $|U|=L-1-2|I|$.  The prime
number theorem gives
\[
\pi(20L\log L)-\pi(L)>2L
\]
for sufficiently large $L$.  After excluding the $|I|$ primes already used
for pairs, more than $2L-|I|>|U|$ primes remain.  Therefore, we may assign a
distinct prime $r_k\in(L,20L\log L]$ to every $k\in U$.

It remains to estimate the product of the selected primes.  Put $r=|I|$ and
$s=|U|=L-1-2r$.  Then
\begin{align*}
\log M
&\leq
r\bigl((2-\alpha_0)\log L+\log2\bigr)
+s\bigl(\log L+\log\log L+\log20\bigr)\\
&=(L-1-\alpha_0r)\log L+s\log\log L+O(L)\\
&\leq
\left(L-1-\frac{\alpha_0L}{100}\right)\log L
+L\log\log L+O(L).
\end{align*}
For sufficiently large $L$, the lower order terms satisfy
\[
L\log\log L+O(L)
\leq\frac{\alpha_0}{200}L\log L,
\]
and therefore
\[
\log M
\leq
\left(1-\frac{\alpha_0}{200}\right)L\log L.
\]
This is the first estimate in the proposition because
$\delta_0=\alpha_0/200$.
Since
\[
\log((L+1)!)=L\log L-L+O(\log L),
\]
we have $\log((L+1)!)\geq L\log L-2L$ for sufficiently large $L$.  Since
$\eta_0=\delta_0/2$, we obtain
\begin{align*}
(1-\eta_0)\log((L+1)!)-\log M
&\geq
(1-\eta_0)(L\log L-2L)-(1-\delta_0)L\log L\\
&=(\delta_0-\eta_0)L\log L-2(1-\eta_0)L>0
\end{align*}
for sufficiently large $L$.  This proves
$M\leq((L+1)!)^{1-\eta_0}$.  The displayed ranges also show that every
selected prime is less than $L!$ once $L$ is sufficiently large.
Choosing $L$ above all the preceding absolute thresholds proves the
proposition.
\end{proof}

\begin{proof}[Proof of Theorem~\ref{thm:covercost}]
Choose the sets and primes from Proposition~\ref{prop:quadratic-cover}.  For
each $i\in I$, the identity \eqref{eq:keyidentity} and the divisibility
$q_i\mid f(i)$ give
\[
i!\equiv(i+2)!\pmod{q_i}.
\]
Hence the edge $e_{q_i,i!}$ covers both $i$ and $i+2$.  For each $k\in U$,
the edge $e_{r_k,k!}$ covers $k$.  All the primes are distinct, so these
edges form a rainbow edge cover $\mathcal C_L$.  Proposition
\ref{prop:quadratic-cover} gives
\[
\kappa(L)\leq w(\mathcal C_L)
=\log M\leq(1-\delta_0)L\log L.
\]
Taking $\delta=\delta_0$ proves the theorem.
\end{proof}


\subsection{Excluding a set of primes}
The proof of Proposition~\ref{prop:quadratic-cover} gives more pairs
than the construction uses. We can therefore discard a prescribed set
of primes and still obtain the same bound. We also record how the
weight depends on the number of chosen pairs.

\begin{theorem}\label{thm:robust}
Let $\alpha_0,\delta_0,\eta_0$ be as in
Proposition~\ref{prop:quadratic-cover}. There is an absolute $L_2$ such
that, for every $L\geq L_2$ and every set $\mathcal F_0$ of at most
$\lfloor L/100\rfloor$ primes, the following holds. For each integer
\[
0\leq r\leq\lceil L/100\rceil,\qquad s=L-1-2r,
\]
there are $r$ disjoint designated pairs with distinct primes
$q_i\in(L,2L^{2-\alpha_0})\setminus\mathcal F_0$ and $s$ distinct
singleton primes in $(L,20L\log L]\setminus\mathcal F_0$, all different
from the pair primes, giving a full cover of modulus $M$ with
\begin{equation}\label{eq:robustcost}
\log M\leq (L-1-\alpha_0r)\log L+s\log\log L+L\log20.
\end{equation}
All selected primes are at most $L^2$ and less than $L!$.
For $r=\lceil L/100\rceil$ one has, uniformly in $\mathcal F_0$,
\[
\log M\leq(1-\delta_0)L\log L,
\qquad M\leq((L+1)!)^{1-\eta_0}.
\]
\end{theorem}
\begin{proof}
The proof of Proposition~\ref{prop:quadratic-cover}, before the final
subfamily is taken, provides at least $9L/400$ pairwise disjoint
pairs with distinct prime colors. Deleting all colors in $\mathcal F_0$
removes at most $L/100$ pairs, leaving at least $L/80$.
This is more than $\lceil L/100\rceil$ for large $L$, so any indicated
number $r$ may be retained.

The prime number theorem gives more than $2L$ primes in
$(L,20L\log L]$ for all sufficiently large $L$. Excluding the forbidden
primes and the $r$ pair colors still leaves more than $s$ choices.
The logarithmic calculation is
\[
\begin{aligned}
\log M
&\leq r((2-\alpha_0)\log L+\log2)
+s(\log L+\log\log L+\log20)\\
&\leq(L-1-\alpha_0r)\log L+s\log\log L+L\log20.
\end{aligned}
\]
For $r=\lceil L/100\rceil$, the last two terms are absorbed exactly as
in Proposition~\ref{prop:quadratic-cover}. Since $\alpha_0>0$ is fixed,
$2L^{2-\alpha_0}\leq L^2$ and $20L\log L\leq L^2<L!$ eventually.
Every threshold depends only on the original absolute constants, not on
the size of the forbidden primes or on their locations.
\end{proof}

\begin{corollary}[Counting the moduli]\label{cor:modulifamily}
Under the hypotheses of Theorem~\ref{thm:robust}, put
$r=\lceil L/100\rceil$, $s=L-1-2r$, and $N=\lfloor L/80\rfloor$.
For all sufficiently large $L$, there are at least
\[
\binom Nr\binom{4L}s\geq2^L
\]
distinct moduli for full covers satisfying its stated power bound and avoiding
$\mathcal F_0$. They support a probability law such that, for every
integer $d\geq1$,
\[
\mathbb P(d\mid M)\leq(9/10)^{\omega(d)}.
\]
\end{corollary}
\begin{proof}
Fix $N$ of the surviving pair colors. The prime number theorem gives
$(20+o(1))L$ primes in $(L,20L\log L]$, so choose a pool of $4L$ primes
avoiding $\mathcal F_0$ and all $N$ pair colors. Select an $r$-subset of
the pair pool and an $s$-subset of the singleton pool, uniformly and
independently. Every choice gives a cover after assigning its singleton
primes to the remaining indices. Unique factorization recovers both
selected subsets from $M$, proving the number of distinct moduli.
Also $\binom{4L}s\geq(4L/s)^s\geq4^s\geq2^L$ eventually.

If $d$ has a repeated prime factor or a factor outside the pools, its
divisibility probability is zero. Otherwise, if it uses $u$ pair primes
and $v$ singleton primes, the probability is
\[
\frac{(r)_u}{(N)_u}\frac{(s)_v}{(4L)_v},
\]
interpreted as zero if too many primes are prescribed.
Here $(a)_j=a(a-1)\cdots(a-j+1)$, and $(a)_0=1$.
Sampling without replacement bounds the ratios by $(r/N)^u$ and
$(s/(4L))^v$. Eventually $r/N\leq9/10$ and $s/(4L)\leq1/4$, proving
the assertion.
\end{proof}

\section{Counting covering classes}\label{sec:ensemble}

Different assignments of indices to primes can give the same residue
class. To count distinct classes, we first bound how often a factorial
residue occurs. The estimates hold for every prime $q>L$.

\subsection{Equal factorial residues}
For $U\subseteq\{1,\ldots,L\}$, put
\[
m_{q,U}(a)=\#\{k\in U:k!\equiv a\pmod q\},
\qquad
\mathcal E_q(U)=\sum_{a\bmod q}m_{q,U}(a)^2.
\]

\begin{proposition}\label{prop:fibresize}
For every prime $q>L$ and $a\pmod q$,
\[
m_{q,\{1,\ldots,L\}}(a)
\leq F(L):=\min_{1\leq h\leq L}
\left(\left\lceil\frac Lh\right\rceil+\frac{h(h-1)}2\right).
\]
In particular,
\[
F(L)\leq\frac32L^{2/3}+O(L^{1/3}).
\]
\end{proposition}
\begin{proof}
Partition $\{1,\ldots,L\}$ into $\lceil L/h\rceil$ consecutive blocks of
length at most $h$. Within each block join consecutive occurrences of
$a$. If $m$ is their total number, this produces at least
$m-\lceil L/h\rceil$ links. A link of length $d<h$ starting at $i$ gives
\[
\prod_{j=1}^d(i+j)-1\equiv0\pmod q.
\]
The polynomial is monic of degree $d$, and the starting indices are
distinct modulo $q$. There are therefore at most $d$ such links, over
all blocks together. Summing over $d$ gives the finite bound.
Taking $h=\lceil L^{1/3}\rceil$ gives the asymptotic estimate.
\end{proof}

\begin{corollary}\label{cor:globallower}
For every sufficiently large $L$,
\[
\kappa(L)\geq\frac{L-1}{F(L)}\log(L+1)
\geq(2/3-o(1))L^{1/3}\log L.
\]
\end{corollary}
\begin{proof}
A cover using $m$ prime colors covers at most $mF(L)$ vertices.
Thus $mF(L)\geq L-1$, and each color costs at least $\log(L+1)$.
\end{proof}

\begin{lemma}\label{lem:energy}
For every $L\geq1$, every prime $q>L$, and every
$U\subseteq\{1,\ldots,L\}$,
\[
\mathcal E_q(U)\leq6L^{3/2}.
\]
\end{lemma}
\begin{proof}
Let $h=\lceil\sqrt L\rceil$ and partition the interval into
$b=\lceil L/h\rceil$ blocks. Write $m_{a,B}$ for the occurrences of
residue $a$ from $U$ in block $B$. Cauchy--Schwarz gives
\[
\mathcal E_q(U)\leq b\sum_{a,B}m_{a,B}^2.
\]
The diagonal terms total $|U|$. For distance $1\leq d<h$, the
polynomial in the previous proof has at most $d$ possible starts, so
there are at most $d$ equal-residue unordered pairs of this distance
within the blocks. Consequently,
\[
\mathcal E_q(U)\leq b\bigl(|U|+h(h-1)\bigr)
\leq(\sqrt L+1)(2L+\sqrt L)\leq6L^{3/2}.
\]
The estimate counts ordered pairs, including equal indices.
\end{proof}

The $L^{3/2}$ order of Lemma~\ref{lem:energy} is also the $\ell=1$
case of \cite[Theorem~1]{GaraevLucaShparlinski}, with their interval
starting at $0$ and containing $L$ indices. The proof above gives the
stated constant uniformly in $q$.

\subsection{Images of permutations}

\begin{lemma}\label{lem:permutationimage}
Let $U$ have $s\geq1$ elements, and let $f_j:U\to\Omega_j$ be maps to
finite sets, for $1\leq j\leq s$. Define
\[
E_j=\#\{(u,v)\in U^2:f_j(u)=f_j(v)\}.
\]
As $(u_1,\ldots,u_s)$ ranges over all permutations of $U$, the number
$N$ of distinct vectors $(f_1(u_1),\ldots,f_s(u_s))$ satisfies
\[
N\geq\frac{s!\,s^s}{\prod_{j=1}^s E_j}.
\]
\end{lemma}
\begin{proof}
For a vector $z$ in the image, let $n_z$ be its number of preimages, and put
$p_z=n_z/s!$. If $m_j(v)=|f_j^{-1}(v)|$, then
$n_z\leq\prod_jm_j(z_j)$, since allowing repeated elements only
increases the number of possible preimages.
For a uniform random permutation, its $j$th entry is uniform on $U$.
Thus
\[
\mathbb E\,m_j(f_j(u_j))=E_j/s.
\]
Concavity of the logarithm gives
\[
\begin{aligned}
\log(s!)
&=\sum_zp_z\log(1/p_z)+\sum_zp_z\log n_z\\
&\leq\log N+\sum_j\mathbb E\log m_j(f_j(u_j))\\
&\leq\log N+\sum_j\log(E_j/s).
\end{aligned}
\]
For the first inequality, Jensen's inequality gives
\[
\sum_zp_z\log(1/p_z)\leq\log\left(\sum_z p_z\frac1{p_z}\right)=\log N.
\]
Rearranging and exponentiating proves the lemma.
\end{proof}

\begin{theorem}[Covering classes for a fixed modulus]
\label{thm:ensemble}
For all sufficiently large $L$, let $\mathcal F_0$ be any set of at most
$\lfloor L/100\rfloor$ primes and let
$0\leq r\leq\lceil L/100\rceil$, $s=L-1-2r$.
There exist a modulus $M$ satisfying \eqref{eq:robustcost} and a set
$\mathcal A\subseteq(\mathbb Z/M\mathbb Z)^\times$ such that every
$a\in\mathcal A$ covers all indices $2\leq k\leq L$ by primes dividing
$M$, no prime in $\mathcal F_0$ divides $M$, and
\begin{equation}\label{eq:ensemblecount}
|\mathcal A|\geq
\frac{s!\,s^s}{(6L^{3/2})^s}
\geq\left(\frac{s^2}{6eL^{3/2}}\right)^s.
\end{equation}
For each fixed $0<\rho\leq1/100$, taking $r=\lceil\rho L\rceil$ gives
\[
M\leq X^{1-\alpha_0\rho/4},\qquad
|\mathcal A|\geq X^{1/2-\rho-o(1)},\qquad X=(L+1)!.
\]
The error term is uniform in $\mathcal F_0$.
\end{theorem}
\begin{proof}
Fix the pairs and the $s$ singleton primes $p_1,\ldots,p_s$ from
Theorem~\ref{thm:robust}. Let $U$ be the indices not in the pairs.
Keep the pair congruences fixed. For each permutation
$(u_1,\ldots,u_s)$ of $U$, impose
\[
a\equiv u_j!\pmod{p_j}\qquad(1\leq j\leq s).
\]
Each system is compatible and reduced, and covers all the remaining
indices. All systems have the same modulus $M$. By CRT, different
vectors of singleton residues give different classes modulo $M$.
Apply Lemma~\ref{lem:permutationimage} to $f_j(u)=u!\bmod p_j$ and
then Lemma~\ref{lem:energy}. The last inequality in
\eqref{eq:ensemblecount} follows from $s!\geq(s/e)^s$.

For fixed $\rho$, equation \eqref{eq:robustcost} gives
$\log M\leq(1-\alpha_0\rho)L\log L+O(L\log\log L)$.
Since $\log X=L\log L-L+O(\log L)$, this implies the asserted bound
for $M$ after increasing the threshold. Also $s=(1-2\rho)L+O(1)$,
so taking logarithms in \eqref{eq:ensemblecount} gives
\[
\log|\mathcal A|\geq(1/2-\rho)L\log L-O_\rho(L).
\]
This proves the fixed-$\rho$ assertions.
\end{proof}

\begin{corollary}\label{cor:halfensemble}
Write $\alpha=\alpha_0$. For every sufficiently large $L$ and every
forbidden set as in Theorem~\ref{thm:ensemble}, there are $M$ and
$\mathcal A$ as in that theorem such that
\[
M\leq Xe^{-5L},\qquad
|\mathcal A|\geq X^{1/2}
\exp\left(-\frac3\alpha L\log\log L\right),\qquad X=(L+1)!.
\]
In particular, $|\mathcal A|\geq X^{1/2-o(1)}$.
\end{corollary}
\begin{proof}
Choose
\[
r=\left\lceil\frac{L(\log\log L+10)}{\alpha\log L}\right\rceil.
\]
Since $\alpha>0$ is fixed, eventually $r\leq\lceil L/100\rceil$ and
$s=L-1-2r\geq L/2$. Equation \eqref{eq:robustcost} yields
\[
\log M\leq L\log L-L(\log\log L+10)+L\log\log L+L\log20
\leq L\log L-7L.
\]
As $\log X\geq L\log L-2L$ eventually, $M\leq Xe^{-5L}$.
Taking logarithms in \eqref{eq:ensemblecount} and using $s\geq L/2$
gives, with an absolute constant $C$,
\[
\begin{aligned}
\log|\mathcal A|
&\geq\tfrac12s\log L-CL\\
&\geq\tfrac12L\log L-
\frac{L}{\alpha}(\log\log L+10)-CL-\tfrac32\log L\\
&\geq\tfrac12\log X-\frac3\alpha L\log\log L
\end{aligned}
\]
for sufficiently large $L$. All estimates are uniform in the forbidden
set. Finally $L\log\log L=o(\log X)$.
\end{proof}
\section{Triples with small span and lower bounds for covers}\label{sec:locality}

The construction above groups indices in pairs. To reduce the weight
further, one may try to use fibers containing three or more indices.
We now bound the number of triples of equal factorial residues whose
first and last indices are close together. We count a triple separately
for each prime for which the equality holds.

For $L\geq3$ and $1\leq D\leq L$, define
\[
T(L,D)=\sum_{\substack{q>L\\q\text{ prime}}}
\#\{2\leq i<j<k\leq L:k-i\leq D,\ i!\equiv j!\equiv k!\pmod q\}.
\]
The sum is finite because each triple imposes divisibility of fixed
nonzero integers. Put $\mathcal B(D)=0$ for $D\leq3$, and otherwise put
\begin{equation}\label{eq:resultantbudget}
\mathcal B(D)=\sum_{b=4}^D
\left(\frac{b(b-3)}2\log2+
\frac{b^3-13b+12}{6}\log b\right).
\end{equation}

\begin{theorem}\label{thm:shorttriples}
For every $L\geq3$ and $1\leq D\leq L$,
\[
T(L,D)\leq\frac{\mathcal B(D)}{\log(L+1)}.
\]
Moreover,
\[
\mathcal B(D)=\frac{D^4}{24}\log D-\frac{D^4}{96}+O(D^3\log D).
\]
Thus for fixed $c>0$,
\[
T(L,\lfloor cL^{1/4}\rfloor)\leq(c^4/96+o(1))L,
\]
and $T(L,D)=o(L)$ whenever $D=o(L^{1/4})$.
\end{theorem}

\subsection{Common roots and resultants}
For $d\geq1$, set
\[
P_d(x)=\prod_{j=1}^d(x+j)-1.
\]
The irreducibility of these polynomials is classical; it is recorded, for
example, in the proof of \cite[Theorem~2.1]{KlurmanMunsch}. We give the
proof together with the resultant bound needed below.

\begin{lemma}\label{lem:resultantheight}
The polynomials $P_d$ are irreducible over $\mathbb Q$. For $1\leq a<b$,
\[
0<|\operatorname{Res}(P_a,P_b)|\leq2^a b^{a(b-a)}.
\]
\end{lemma}
\begin{proof}
If $P_d=FG$ were a nontrivial factorization, Gauss's lemma would give
monic integral factors, with $1\leq m:=\deg F\leq d/2$ after
interchanging them. For $j=1,\ldots,d$, $F(-j)G(-j)=-1$, hence
$F(-j)=\pm1$. Thus $F^2-1$ has $d$ distinct roots. If $2m<d$ this is
impossible. If $2m=d$, monicity gives
\[
F^2-1=\prod_{j=1}^d(x+j)=P_d+1,
\]
so $P_d=F^2-2$. The divisibility $F\mid P_d$ would imply $F\mid2$,
again impossible.

The distinct degrees of the irreducible polynomials imply that their
resultant is a nonzero integer. If $P_a(\zeta)=0$, then
$\prod_{j=1}^a(\zeta+j)=1$, so $|\zeta+j_0|\leq1$ for some
$j_0\leq a$. For $a<j\leq b$ this gives $|\zeta+j|\leq b$.
Consequently,
\[
|P_b(\zeta)|=
\left|\prod_{j=a+1}^b(\zeta+j)-1\right|
\leq2b^{b-a}.
\]
The product of these bounds over the $a$ roots of $P_a$ is the asserted
resultant bound.
\end{proof}

\begin{lemma}\label{lem:rootvaluation}
If $F,G\in\mathbb Z[x]$ are monic nonconstant polynomials with nonzero
resultant, then the number of their
distinct common roots modulo a prime $q$ is at most
$v_q(\operatorname{Res}(F,G))$.
\end{lemma}
\begin{proof}
Let $m=\deg F$, $n=\deg G$, and consider the Sylvester map
$(U,V)\mapsto UF+VG$, where $\deg U<n$ and $\deg V<m$.
Monicity implies that the reductions have the same degrees, so the
number $d$ of their distinct common roots is at most $\min(m,n)$.
Every polynomial in the image vanishes at these roots. Evaluation at the $d$ points imposes
$d$ independent linear conditions on polynomials of degree less than
$m+n$. Thus the integer Sylvester matrix has rank deficiency at least
$d$ modulo $q$.

Lift invertible row operations over $\mathbb F_q$ to an integral matrix
whose determinant is not divisible by $q$. At least $d$ rows of the
resulting matrix are divisible by $q$, so its determinant is divisible
by $q^d$. Its determinant differs from the original resultant by a
$q$-adic unit. Hence $q^d\mid\operatorname{Res}(F,G)$.
\end{proof}

\begin{proof}[Proof of Theorem~\ref{thm:shorttriples}]
Put $a=j-i$ and $b=k-i$. By Proposition~\ref{prop:criterion}, a triple
collision gives a common root $i$ of $P_a$ and $P_b$ modulo $q$.
For fixed $a,b,q$, the possible starts $2\leq i\leq L-b$ are distinct
modulo $q>L$. Thus Lemma~\ref{lem:rootvaluation} bounds their number
by $v_q(R_{a,b})$, where $R_{a,b}=\operatorname{Res}(P_a,P_b)$.
Summing over $q>L$ gives
\[
\begin{aligned}
\log(L+1)\sum_{q>L}\#\{\text{valid starts for }a,b,q\}
&\leq\sum_{q>L}v_q(R_{a,b})\log q\\
&\leq\log|R_{a,b}|\\
&\leq a\log2+a(b-a)\log b.
\end{aligned}
\]
Adjacent factorials do not collide in the index range: cancelling $i!$
from $(i+1)!\equiv i!\pmod q$ would give $q\mid i$.
Therefore $2\leq a\leq b-2$ and $b\geq4$. Summing the last bound over
these gaps yields \eqref{eq:resultantbudget}, since
\[
\sum_{a=2}^{b-2}a=\frac{b(b-3)}2,
\qquad
\sum_{a=2}^{b-2}a(b-a)=\frac{b^3-13b+12}{6}.
\]
Finally, comparison with the integral gives
\[
\sum_{b\leq D}b^3\log b
=\frac{D^4}{4}\log D-\frac{D^4}{16}+O(D^3\log D).
\]
The remaining sums in \eqref{eq:resultantbudget} are of lower order.
This proves the expansion and its consequences. For bounded $D$, the
$o(L)$ assertion follows directly from the finite bound.
\end{proof}

\subsection{Lower bounds for covers}
We call a choice of one nonempty set
$B_q\subseteq E_{q,a_q}$ for each selected prime, covering the indicated
vertices, a \emph{designated cover}. These sets may overlap. A consecutive
triple consists of three successive elements in the increasing
enumeration of one $B_q$; the elements need not be consecutive integers.

\begin{theorem}\label{thm:localcover}
Let a designated cover have weight $w$, and let
$n=|\bigcup_q B_q|$. Let $K_{>D}$ count its prime-labelled consecutive
triples of span greater than $D$. Then
\[
K_{>D}\geq n-\frac{2w+\mathcal B(D)}{\log(L+1)}.
\]
In particular, if all its consecutive triples have span at most $D$,
then
\[
w\geq\tfrac12\bigl(n\log(L+1)-\mathcal B(D)\bigr).
\]
\end{theorem}
\begin{proof}
Write $m$ for the number of selected primes and $b_q=|B_q|$.
There are $\sum_q(b_q-2)_+\geq\sum_qb_q-2m\geq n-2m$
consecutive triples. At most $T(L,D)$ of them have span at most $D$.
They are distinct as prime-labelled triples, even when different colors
have overlapping designated sets. Now apply Theorem~\ref{thm:shorttriples}
and $m\leq w/\log(L+1)$.
\end{proof}

\begin{corollary}\label{cor:localcost}
Suppose $c>0$ and $0\leq u<1$ are fixed, and a sequence of designated
covers covers at least $(1-u)(L-1)$ indices. If every consecutive-triple span
is at most $cL^{1/4}$, then
\[
\frac{w}{\log((L+1)!)}\geq\frac{1-u}{2}-\frac{c^4}{192}-o(1).
\]
For full covers this gives $1/2-c^4/192-o(1)$. If every span is
$o(L^{1/4})$, the full-cover lower bound is $1/2-o(1)$.
\end{corollary}
\begin{proof}
Use $D=\lfloor cL^{1/4}\rfloor$ in Theorem~\ref{thm:localcover} and
$\log((L+1)!)=L\log L-L+O(\log L)$. For a full cover, use $n=L-1$
directly. For a varying $D=o(L^{1/4})$, use $\mathcal B(D)/\log L=o(L)$.
\end{proof}

For example, if a full cover has $w\leq\theta\log((L+1)!)$ for a fixed
$\theta<1/2$, then for each fixed $c>0$ it has at least
\[
(1-2\theta-c^4/96-o(1))L
\]
consecutive triples of span exceeding $cL^{1/4}$, counted with their
primes, whenever the coefficient is positive.

Suppose a designated set is connected by a graph whose edges join
indices at distance at most $h$. Every consecutive gap in that set is
then at most $h$, since a larger gap would disconnect the graph. Thus its
consecutive-triple spans are at most $2h$, and the preceding inequalities
apply with $D=\lfloor2h\rfloor$.
\section{Rough integers in prescribed progressions}\label{sec:arithmetic}

We now count rough integers in the residue classes given by a cover.
The Chinese remainder theorem and inclusion--exclusion give the
following estimate, which also allows a prescribed congruence.

\begin{proposition}[Counting rough integers]\label{prop:counting}
Let $L\geq2$, $W=\prod_{p\leq L}p$, and let $M,d\geq1$ satisfy
$(M,dW)=1$. Suppose $(A,M)=(b,d)=1$. Put
\[
W_d=\prod_{p\leq L,\ p\nmid d}p,
\qquad V_d(L)=\prod_{p\leq L,\ p\nmid d}(1-1/p).
\]
For real $u<v$, the number $N_{A,b}(u,v)$ of integers $n\in(u,v]$ with
\[
n\equiv A\pmod M,\qquad n\equiv b\pmod d,
\qquad(n,W)=1
\]
satisfies
\begin{equation}\label{eq:exactcount}
\left|N_{A,b}(u,v)-\frac{v-u}{Md}V_d(L)\right|
\leq2^{\omega(W_d)}.
\end{equation}
The bound is uniform in every displayed parameter.
\end{proposition}
\begin{proof}
The two prescribed congruences give a unique class modulo $Md$.
Primes dividing $d$ cannot divide $n$ because $(b,d)=1$.
For each $e\mid W_d$, the additional condition $e\mid n$ gives a unique
class modulo $Mde$, since $(e,Md)=1$. Its count in $(u,v]$ differs from
$(v-u)/(Mde)$ by at most one. Inclusion--exclusion therefore gives
\[
N_{A,b}(u,v)=\sum_{e\mid W_d}\mu(e)
\left(\frac{v-u}{Mde}+\epsilon_e\right),\qquad |\epsilon_e|\leq1.
\]
The main term factors as $\frac{v-u}{Md}\prod_{p\mid W_d}(1-1/p)$,
and there are $2^{\omega(W_d)}$ error terms.
\end{proof}

Let $\mathcal A$ be a set of distinct classes modulo a common modulus
$M$, each giving a full cover of $2!,\ldots,L!$ with primes below $L!$.
Assume $(M,d)=1$ and $(b,d)=1$. For $X=(L+1)!$, formula
\eqref{eq:exactcount} gives at least
\begin{equation}\label{eq:ensembletransference}
|\mathcal A|\left(
\frac{X-2L!}{Md}V_d(L)-2^{\omega(W_d)}
\right)
\end{equation}
$L$-rough integers in these classes and in $b\pmod d$.
The classes are disjoint, because equality of the combined classes
modulo $Md$ implies equality modulo $M$.

All these integers satisfy the factorial-difference conditions. For
$2\leq k\leq L$, some covering prime satisfies
\[
q\mid n-k!,\qquad 1<q<L!<n-k!.
\]
Thus it is a proper divisor. Roughness makes $n$ odd, and $n-1>2$,
so the case $k=1$ follows as well.

\subsection{Counting in one covering class}

\begin{corollary}\label{cor:costcount}
Let $G_L$ count the auxiliary solutions in $(2L!,X]$, $X=(L+1)!$.
For every sufficiently large $L$,
\[
G_L\geq (X-2L!)e^{-\kappa(L)}
\prod_{p\leq L}(1-1/p)-2^{\pi(L)}.
\]
More generally, for any fixed $\lambda>\log2$, uniformly over full
covers with
\[
M\leq X\exp\left(-\lambda\frac{L}{\log L}\right),
\]
the number of $L$-rough integers in their single CRT class and in
$(2L!,X]$ is
\[
(e^{-\gamma}+o(1))\frac{X-2L!}{M\log L},
\]
where $\gamma$ is Euler's constant. Every integer counted is an auxiliary
solution. This is an asymptotic for the selected progression, not for
all auxiliary solutions.
\end{corollary}
\begin{proof}
Take $d=1$ and a minimum-cost cover in
Proposition~\ref{prop:counting} for the first assertion.
For the second, Mertens's product estimate gives
$V_1(L)\sim e^{-\gamma}/\log L$, and the prime number theorem gives
$2^{\pi(L)}=\exp((\log2+o(1))L/\log L)$.
The main term in \eqref{eq:exactcount} divided by its error bound tends
to infinity uniformly under the stated modulus condition, because
$\lambda>\log2$. This proves the asymptotic. The classical estimates are
\cite[Theorems~3.4(c) and~8.1]{Koukoulopoulos}.
\end{proof}

\begin{theorem}\label{thm:transference}
Fix $0<\eta<1$. For every sufficiently large $L$, depending only on
$\eta$, a full cover with $M\leq X^{1-\eta}$, $X=(L+1)!$, yields at
least $X^\eta/(100\log L)$ auxiliary solutions in $(2L!,X]$.
\end{theorem}
\begin{proof}
Use \eqref{eq:exactcount} with $d=1$ and the cover's CRT class.
For sufficiently large $L$, $X-2L!\geq X/2$ and
$V_1(L)\geq1/(4\log L)$. Moreover,
\[
2^{\pi(L)}=\exp(O(L/\log L))
=o\left(\frac{X^\eta}{\log L}\right).
\]
The count is therefore at least $X/(16M\log L)$ for sufficiently
large $L$, which exceeds the asserted bound. The argument preceding
Corollary~\ref{cor:costcount} shows that all the required differences
are composite.
\end{proof}

\begin{proof}[Proof of Theorem~\ref{thm:main}]
Decrease the constant $\delta$ in Theorem~\ref{thm:covercost} if
necessary so that $0<\delta<1$, and put $\eta=\delta/2$.
The minimum-cost cover exists for sufficiently large $L$ by
Proposition~\ref{prop:baseline}. Since
\[
\log M=\kappa(L)\leq(1-\delta)L\log L,
\qquad \log X=L\log L-L+O(\log L),
\]
one has $M\leq X^{1-\eta}$ eventually.
Theorem~\ref{thm:transference} proves the assertion.
\end{proof}

\subsection{Counting in a prescribed residue class}

\begin{proof}[Proof of Theorem~\ref{thm:abundance}]
Write $\alpha=\alpha_0$. For a given $d\leq e^L$, let $\mathcal F_0$
be the set of prime divisors of $d$ exceeding $L$. Then
\[
|\mathcal F_0|\log(L+1)\leq\log d\leq L,
\]
so $|\mathcal F_0|\leq\lfloor L/100\rfloor$ for all sufficiently large
$L$, uniformly in $d$. Use Corollary~\ref{cor:halfensemble} avoiding this
set. Its modulus satisfies
\[
(M,d)=1,\qquad M\leq Xe^{-5L},\qquad
|\mathcal A|\geq X^{1/2}
\exp\left(-\frac3\alpha L\log\log L\right).
\]
All selected primes exceed $L$, so $(M,W)=1$ as well.

For any $(b,d)=1$, apply \eqref{eq:ensembletransference}.
Since $d\leq e^L$, $V_d(L)\geq V_1(L)\geq1/(4\log L)$, and
$2^{\omega(W_d)}\leq2^L$, the number of rough integers contributed by
each combined class is at least
\[
\frac{e^{4L}}{8\log L}-2^L\geq1
\]
for all sufficiently large $L$. Hence the total is at least
$|\mathcal A|$. The preceding argument gives compositeness of the
factorial differences. Finally,
\[
\frac{(3/\alpha)L\log\log L}{\log X}\longrightarrow0.
\]
For each $\varepsilon>0$, increasing the threshold makes the lower
bound at least $X^{1/2-\varepsilon}$. Every threshold is independent of
$d$ and $b$, as required.
\end{proof}
\section{Computational examples and further questions}\label{sec:conclusion}

\subsection{Examples of covers}
The supplementary files contain full covers at three finite values of
$L$. Each selected prime, residue, complete fiber, and designated subset
is listed. The accompanying program recomputes the residues, checks
primality by trial division, and verifies that the primes are distinct
and the sets cover all indices. It checks the bounds below by comparing
integer powers.

\begin{center}
\small
\begin{tabular}{rrrrl}
\toprule
$L$ & Primes & Consecutive triples & Minimum span & Bound\\
\midrule
300 & 99 & 106 & 24 & $M^{100}<X^{46}$\\
1000 & 296 & 411 & 39 & $M^{100}<X^{39}$\\
3000 & 810 & 1383 & 30 & $M^{50}<X^{17}$\\
\bottomrule
\end{tabular}
\end{center}
Here $X=(L+1)!$, and the triples are consecutive triples of the designated
subsets. In all three certificates every prime lies in $(L,L^2]$.
No listed triple has span at most $\lfloor L^{1/4}\rfloor$.
As one concrete example from the $L=300$ certificate,
\[
17!\equiv81!\equiv87!\equiv176!\equiv198!\equiv293!
\equiv70\pmod{337}.
\]
These examples give the displayed bounds for $\kappa(L)$ at the three
specified values. They give no asymptotic bound. The supplementary
programs also check the inequality in Lemma~\ref{lem:permutationimage},
the progression counts, and resultants for small degrees. These
computations are separate from the proofs.

\subsection{Further questions}
The upper and lower bounds for the minimum weight are far apart:
\[
(2/3-o(1))L^{1/3}\log L
\leq\kappa(L)\leq(1-\delta)L\log L.
\]
Our arguments do not determine its order or prove
$\kappa(L)=o(L\log L)$. Corollary~\ref{cor:localcost} gives one
restriction on improving the upper bound: a cover whose weight is less
than half the leading weight of the individual assignments, by a fixed
proportion, must contain a positive proportion of triples with large
span. Choosing such triples with distinct primes and disjoint indices
is a further problem.

Proposition~\ref{prop:completion} gives a way to use additional equalities
between factorial residues. A compatible family of designated subsets
with a large total saving $\sum_F(|B_F|\log T-\log q_F)$ gives a cover
of small weight, provided enough unused primes of size at most $T$ are
available for the remaining indices. The difficulty is to find such
subsets and choose them without repeating primes or indices. A lower
bound for the number of distinct factorial residues alone does not give
this choice.

One can also ask how many covering classes a fixed modulus admits.
Lemma~\ref{lem:permutationimage} bounds this number using counts of
equal factorial residues at each prime. Stronger estimates for these
counts may improve the $X^{1/2-o(1)}$ lower bound even if the bound for
$\kappa(L)$ stays the same. The integers counted by this method are
rough integers; primality is a separate condition.

\section*{Acknowledgments}

This work was carried out as part of the 2026 Research in Industrial Projects
for Students program at the Institute for Pure \& Applied Mathematics,
UCLA.  The author thanks IPAM for providing a collaborative research
environment and acknowledges OpenAI, the industrial sponsor of the project,
for its support and for providing access to the computational tools used
during the project.

\section*{Use of AI tools}

The mathematical exploration and initial proof development used OpenAI's
ChatGPT and Codex extensively. These tools also assisted with
computation, literature searches, algebraic checks, and manuscript
preparation. The author selected and reframed the problem, proposed the
broader hypergraph interpretation, and directed the successive revisions.
Aristotle (Harmonic) produced the Lean formalization of the original
auxiliary theorem. The extent of that formalization is stated in the
introduction.

\begin{thebibliography}{99}

\bibitem{BanksLucaShparlinskiStichtenoth}
W.~D. Banks, F.~Luca, I.~E. Shparlinski, and H.~Stichtenoth,
\emph{On the value set of $n!$ modulo a prime},
Turkish J. Math. \textbf{29} (2005), no.~2, 169--174.

\bibitem{Bloom}
T.~F. Bloom,
\emph{Erd\H{o}s Problem 1059},
\href{https://www.erdosproblems.com/1059}
{erdosproblems.com/1059}, accessed August 9, 2026.

\bibitem{GaraevLucaShparlinski}
M.~Z. Garaev, F.~Luca, I.~E. Shparlinski,
\emph{Exponential sums and congruences with factorials},
J. Reine Angew. Math. \textbf{584} (2005), 29--44.

\bibitem{Guy}
R.~K. Guy,
\emph{Unsolved Problems in Number Theory},
3rd ed., Problem Books in Mathematics, Springer, New York, 2004,
Problem A2, pp.~10--12.

\bibitem{Hu}
X.~Hu,
\emph{Factorial residues modulo a prime: beyond the square-root bound},
\href{https://arxiv.org/abs/2608.01781}{arXiv:2608.01781} (2026).

\bibitem{KlurmanMunsch}
O.~Klurman and M.~Munsch,
\emph{Distribution of factorials modulo $p$},
J. Th\'eor. Nombres Bordeaux \textbf{29} (2017), no.~1, 169--177.

\bibitem{Koukoulopoulos}
D.~Koukoulopoulos,
\emph{The Distribution of Prime Numbers},
Graduate Studies in Mathematics, vol.~203,
American Mathematical Society, Providence, RI, 2019.

\bibitem{Land}
J.~Land,
\emph{A proposed solution to Erd\H{o}s Problem 1059: prime avoidance for a sparse set of shifts},
preprint, September 5, 2026,
\href{https://github.com/beetree/math_erdos_1059}
{github.com/beetree/math\_erdos\_1059}.

\bibitem{Palomar}
Palomar Registry,
\emph{Factorial-residue hypergraphs and an auxiliary problem of Erd\H{o}s},
entry PALOMAR-2026-08-27-000003 v1,
\href{https://palomar-registry.org/entry.html?id=PALOMAR-2026-08-27-000003&version=1}
{palomar-registry.org}, 2026.

\end{thebibliography}

\end{document}
